% Semester 1 capstone 1: fast multiplication.
\documentclass{ca-capstone}
\capstonerefs{knuth-taocp2,crandall-pomerance,welschenbach,sagemath}
\begin{document}
% Generated by `make bench CAPSTONE=fast-multiplication` (taskgen.capstone report) from
% bench/results.json: reference time x slack on the reference machine.
% Do not edit: change bench/config.yaml and re-run the benchmark.
\newcommand{\ReqSlack}{10}
\newcommand{\ReqBitsMax}{2048}
\newcommand{\ReqSeconds}{0{,}05}

\capstonetitle{Быстрое умножение больших чисел}

\goal
Требуется реализовать и сравнить несколько алгоритмов умножения целых чисел произвольной
длины, включая умножение через быстрое преобразование Фурье в автоматически подобранном
конечном поле.

\functionality
\begin{enumerate}
  \item Умножение школьным методом и методом Карацубы.
  \item Умножение методом Тоома-Кука в одном фиксированном варианте: Тоом-3.
  \item Прямая реализация линейной и циклической свёртки конечных последовательностей.
  \item Умножение с помощью БПФ в конечном поле:
    \begin{enumerate}[label=\alph*.]
      \item вычисление требуемой длины преобразования и верхней оценки коэффициентов свёртки;
      \item автоматический выбор подходящего простого~$p$, корня из единицы требуемого
        порядка и обратных элементов;
      \item прямое и обратное преобразования, поточечное умножение и точное восстановление
        целочисленного результата.
    \end{enumerate}
  \item Корректная обработка нуля, знаков, чисел существенно разной длины, переносов и
    ошибочных внутренних параметров преобразования.
  \item Сравнение времени работы всех четырёх методов на числах нескольких размеров.
\end{enumerate}

\requirements
\begin{enumerate}
  \item Все методы должны возвращать точное произведение.
  \item Для БПФ должны быть явно обоснованы достаточность выбранного~$p$, порядок корня и
    обратимость преобразования.
  \item Тестовые примеры с входными числами длиной до \textbf{\ReqBitsMax~бит} должны
    выполняться не более чем за \textbf{\ReqSeconds~с} на операцию.
\end{enumerate}

Корректность проверяется сравнением с библиотечным умножением и прямым вычислением свёртки.
Отдельно необходимо проверить корень из единицы, обратное преобразование, переходы оценки коэффициентов
через очередное подходящее простое.

\details
Разрешается использовать любой язык программирования, в частности Python, систему 
компьютерной алгебры SageMath и встроенные большие целые для
вспомогательных операций сложения, сдвигов и модульной арифметики. Реализуемые методы
умножения и БПФ должны быть написаны самостоятельно; встроенное умножение и библиотечные
преобразования можно использовать как независимый оракул и для поэлементных операций с
машинными цифрами. При поиске~$p$ разрешается использовать библиотечную функцию проверки простоты.

Не требуются вещественное БПФ, отдельные реализации ТЧП Ферма и Мерсенна, восстановление
по нескольким модулям и автоматический выбор семейства алгоритма по размеру входа.

Сохраняйте параметры преобразования $(N,p,\omega)$ и оценку коэффициентов как промежуточные
результаты. Для остальных методов желательно сохранять глубину рекурсии и промежуточные результаты.

\schedule
Работа выполняется индивидуально. Срок выполнения -- \textbf{месяц}; ожидаемая трудоёмкость
-- примерно 36-40 часов. Реализация должна быть разработана и написана вами.

\remarks
\begin{enumerate}
  \item Обязательны unit-тесты и сравнение с независимым оракулом.
  \item Оптимизация кода не требуется; важнее корректность параметров и точность
    восстановления результата.
  \item Дополнительные алгоритмы и необязательные расширения в работу не входят.
\end{enumerate}

\references
\end{document}
